\chapter{Results}\label{chapter:results}
\markboth{\textit{CHAPTER 4. Results}}{}

\clearpage
\FloatBarrier

This chapter presents the findings of the thesis in four main sections: (1) characterisation of the transcriptomic atlas and sample distribution; (2) evaluation of model performance across tissues; (3) analysis of developmental gene signatures and their functional annotation; and (4) cross-species validation with human skeletal muscle.


\section{Data Landscape and Sample Characterisation}\label{sec:results:data}

\subsection{Sample Distribution Across Tissues and Developmental Stages}

The final dataset comprised 1,924 samples spanning five tissues and up to five developmental stages. The distribution of samples across tissues and stages is presented in \figref{fig:sample_distribution}.

\begin{figure}[H]
\centering
\includegraphics[width=\textwidth]{fig1_study_design.pdf}
\caption{\textbf{Study design and data landscape.}
\textbf{(a)} Developmental stage definitions mapping chronological age to five ordered stages: infant (0--20 days), early childhood (21--59 days), pre-pubertal (60--149 days), post-pubertal (150--365 days), and adult (>365 days).
\textbf{(b)} Machine learning workflow diagram illustrating the pipeline from data preprocessing to model calibration and evaluation.
\textbf{(c)} Sample distribution heatmap showing tissue-by-stage availability across 1,924 samples from five tissues.
\textbf{(d)} Classification scheme assignments based on per-class sample availability.}
\label{fig:sample_distribution}
\end{figure}

Muscle tissue contributed the largest number of samples ($n = 761$), followed by liver ($n = 309$), lung ($n = 129$), blood ($n = 78$), and brain ($n = 43$). This uneven distribution reflects the relative availability of these tissues in the PigGTEx resource \parencite{teng_compendium_2024, chen_construction_2025, long_comprehensive_2024} and the differential representation of developmental stages across tissues.

\subsection{Class Distribution and Imbalance Analysis}

The distribution of samples across developmental stages varied substantially between tissues. Muscle and liver showed relatively balanced representation across four stages (infant, early childhood, pre-pubertal, and post-pubertal/adult combined). In contrast, brain, blood, and lung samples were concentrated in younger developmental stages, with limited representation of adult animals.

The quantitative distribution was as follows:

\textbf{Muscle} (761 samples) was predominantly composed of post-pubertal/adult samples (401, 52.7\%), followed by infant (194, 25.5\%), pre-pubertal (90, 11.8\%), and early childhood (76, 10.0\%) stages.
\textbf{Liver} (309 samples) showed a more even distribution, with infant (87, 28.2\%), pre-pubertal (83, 26.9\%), and post-pubertal/adult (81, 26.2\%) stages similarly represented, while early childhood (58, 18.8\%) was slightly less common.
\textbf{Brain} (43 samples) was heavily skewed towards young animals, with 28 samples (65.1\%) in the pre-pubertal stage and only 15 samples (34.9\%) in the post-pubertal stage.
\textbf{Blood} (78 samples) and \textbf{Lung} (129 samples) distinctively favored younger stages, with 60.3\% and 79.8\% of samples classified as young, respectively. This contrasted with the mature populations comprising 39.7\% and 20.2\% of the cohorts.

\subsection{Classification Scheme Selection}

Based on the sample distribution, the adaptive scheme selection algorithm assigned classification schemes as follows:

Specifically, muscle and liver were assigned 4-class schemes (infant, early childhood, pre-pubertal, post-pubertal/adult), enabling finer-grained staging. Conversely, brain, blood, and lung were assigned 2-class schemes (young vs. mature), ensuring robust classification despite smaller sample sizes.


\section{Model Performance Assessment}\label{sec:results:performance}

\subsection{Overall Performance Summary}

The machine learning framework achieved robust performance across all five tissues, with mean balanced accuracy of 89.5\% evaluated via nested cross-validation (pooled held-out predictions). Performance varied according to tissue complexity, sample size, and classification scheme granularity.

The performance metrics for each tissue are summarised in \tabref{tab:performance} and visualised in \figref{fig:performance}.

\begin{table}[H]
\centering
\caption{\textbf{Model performance metrics across tissues.} Performance evaluated via nested cross-validation (pooled held-out predictions). 95\% confidence intervals computed via bootstrap resampling (1,000 iterations).}
\label{tab:performance}
\small
\begin{tabular}{lccccc}
\toprule
\textbf{Tissue} & \textbf{Scheme} & \textbf{$n$} & \textbf{Balanced Acc.} & \textbf{F1-macro} & \textbf{Spearman $\rho$} \\
\midrule
Blood & 2-class & 78 & 0.959 & 0.960 & -- \\
Liver & 4-class & 309 & 0.921 & 0.918 & 0.949 \\
Muscle & 4-class & 761 & 0.894 & 0.899 & 0.965 \\
Lung & 2-class & 129 & 0.892 & 0.876 & -- \\
Brain & 2-class & 43 & 0.811 & 0.801 & 0.942 \\
\midrule
\textbf{Mean} & -- & -- & \textbf{0.895} & \textbf{0.891} & -- \\
\bottomrule
\end{tabular}
\end{table}

\begin{figure}[H]
\centering
\includegraphics[width=\textwidth]{fig2_model_performance.pdf}
\caption{\textbf{Classifier performance diagnostics.}
\textbf{(a)} Performance metrics heatmap showing balanced accuracy, macro F1, mean absolute error, and Spearman correlations with 95\% bootstrap confidence intervals across tissues.
\textbf{(b)} Confusion matrices for each tissue showing classification accuracy and error patterns.
\textbf{(c)} Scatter plots of predicted developmental stage versus chronological age, demonstrating strong age--stage correlations.}
\label{fig:performance}
\end{figure}

\subsection{Tissue-Specific Performance}

\subsubsection{Blood: Highest Performance}

Blood achieved the highest balanced accuracy (95.9\%) using the 2-class (pre-pubertal vs.\ post-pubertal) scheme. The macro F1-score of 0.960 indicates excellent precision and recall across both classes, demonstrating robust discrimination between pre-pubertal and post-pubertal animals even with a modest sample size ($n = 78$).

\subsubsection{Liver: Strong 4-Class Classification}

Liver achieved a balanced accuracy of 92.1\%, reflecting its relatively homogeneous cellular composition and balanced class representation. The macro F1-score of 0.918 indicates consistently high precision and recall across all four developmental stages.

The strong Spearman correlation ($\rho = 0.949$) between predicted stage and chronological age demonstrates excellent ordinal consistency. The model correctly ordered the vast majority of samples by developmental stage, with errors primarily occurring at stage boundaries.

\subsubsection{Lung: Strong Binary Classification}

Lung achieved a balanced accuracy of 89.2\% using the 2-class scheme, despite pronounced class imbalance (79.8\% pre-pubertal samples). The macro F1-score of 0.876 confirms effective classification across both classes. The use of class weighting during training enabled the model to learn from the minority class effectively.

\subsubsection{Muscle: Robust 4-Class Classification}

Muscle achieved a balanced accuracy of 89.4\% for 4-class classification, demonstrating that the framework can discriminate multiple developmental stages in this metabolically active tissue. The Spearman correlation of 0.965 indicates near-perfect ordinal consistency.

The large sample size for muscle ($n = 761$) provides statistical power that enables precise confidence interval estimation and reliable identification of developmental markers.

\subsubsection{Brain: Lowest Performance}

The brain showed the lowest balanced accuracy (81.1\%), reflecting the challenges of this tissue. The limited sample size ($n = 43$) constrained the model's ability to learn robust developmental signatures. Additionally, the brain's extreme cellular heterogeneity, comprising diverse neuronal and glial populations \parencite{jiang_cellular_2016}, may obscure developmental signals in bulk RNA-seq data.

Despite the lower overall accuracy, the Spearman correlation of 0.942 indicates that the model correctly orders samples by developmental stage, even when exact classification is incorrect. Notably, the brain's performance (81.1\%) is substantially improved compared to earlier iterations and remains robust for a tissue with only 43 samples.

\subsection{Per-Class Performance Analysis}

Per-class metrics reveal the stages that are easiest and hardest to classify. The per-class precision, recall, and F1-scores are presented in \tabref{tab:perclass}.

\begin{table}[H]
\centering
\caption{\textbf{Per-class performance metrics by tissue.} Precision, recall, and F1 scores for each developmental stage within each tissue.}
\label{tab:perclass}
\small
\begin{tabular}{llcccc}
\toprule
\textbf{Tissue} & \textbf{Stage} & \textbf{$n$} & \textbf{Precision} & \textbf{Recall} & \textbf{F1} \\
\midrule
Muscle & infant & 194 & 0.939 & 0.959 & 0.949 \\
Muscle & early childhood & 76 & 0.838 & 0.816 & 0.827 \\
Muscle & pre-pubertal & 90 & 0.871 & 0.822 & 0.846 \\
Muscle & post-pubertal/adult & 401 & 0.973 & 0.980 & 0.976 \\
\midrule
Brain & pre-pubertal ($<$150d) & 28 & 0.885 & 0.821 & 0.852 \\
Brain & post-pubertal ($\geq$150d) & 15 & 0.706 & 0.800 & 0.750 \\
\midrule
Liver & infant & 87 & 0.965 & 0.954 & 0.960 \\
Liver & early childhood & 58 & 0.887 & 0.948 & 0.917 \\
Liver & pre-pubertal & 83 & 0.910 & 0.855 & 0.882 \\
Liver & post-pubertal/adult & 81 & 0.904 & 0.926 & 0.915 \\
\midrule
Blood & pre-pubertal ($<$150d) & 45 & 0.957 & 0.978 & 0.967 \\
Blood & post-pubertal ($\geq$150d) & 33 & 0.969 & 0.939 & 0.954 \\
\midrule
Lung & pre-pubertal ($<$150d) & 102 & 0.960 & 0.931 & 0.945 \\
Lung & post-pubertal ($\geq$150d) & 27 & 0.767 & 0.852 & 0.807 \\
\bottomrule
\end{tabular}
\end{table}

Several patterns emerged from the per-class analysis. First, extreme stages (infant and post-pubertal/adult) consistently showed higher F1-scores than intermediate stages, reflecting the greater transcriptomic distinctiveness of developmental extremes. Second, early childhood proved to be the most challenging stage to classify in both muscle (F1 = 0.827) and liver (F1 = 0.917), likely due to the transitional nature of the weaning period \parencite{pluske_gastrointestinal_2018}. Finally, the post-pubertal class in brain (F1 = 0.750) and lung (F1 = 0.807) showed lower performance, illustrating the difficulty of learning from limited examples in the minority class.

\subsection{Ordinal Consistency Analysis}

The ordinal nature of developmental staging means that not all errors are equally severe. Misclassifying an infant as early childhood is a smaller error than misclassifying an infant as adult. The confusion matrices in \figref{fig:performance}b reveal the pattern of errors.

For liver, the confusion matrix shows that errors are concentrated on the diagonal and adjacent cells. Only four samples (out of 309 pooled held-out predictions) were misclassified by more than one stage. This near-perfect ordinal consistency reflects the model's ability to capture the continuous nature of development.

For muscle, a similar pattern emerges. The majority of errors involve adjacent stages, with very few distant misclassifications. The early childhood $\leftrightarrow$ pre-pubertal boundary shows the highest error rate, consistent with the biological similarity of these transitional stages.

For the 2-class tissues (brain, blood, lung), ordinal analysis is less informative, as there is only one possible misclassification type. However, the binary classification still reflects the model's ability to distinguish young from mature animals.

\subsection{Relationship Between Predicted Stage and Chronological Age}

The scatter plots in \figref{fig:performance}c demonstrate the relationship between model-predicted developmental stage and chronological age. For tissues with continuous age data, strong positive correlations were observed: muscle ($\rho = 0.965$), liver ($\rho = 0.949$), and brain ($\rho = 0.942$).

These correlations confirm that the model's predictions align with chronological age, as expected given the age-based stage definitions. However, the imperfect correlation reflects inter-individual variability in developmental rate and the limitations of using chronological age as ground truth.


\section{Feature Analysis and Biological Interpretation}\label{sec:results:features}

\subsection{Top Developmental Markers Per Tissue}

Feature importance analysis identified the genes most informative for developmental stage classification in each tissue. The top markers reflect tissue-specific developmental programs.

\subsubsection{Muscle Developmental Markers}

The top developmental markers for muscle include genes involved in contractile function, neuromuscular signalling, and metabolic adaptation:

Key genes included \textbf{FGFRL1} (Fibroblast growth factor receptor-like 1), a decoy receptor modulating FGF signalling \parencite{amann_fgfrl1_2014}, and \textbf{SORCS2} (Sortilin-related VPS10 domain-containing receptor 2), which is involved in neurotrophin signalling \parencite{thomasen_sorcs2_2023}. Other significant markers were \textbf{ACHE} (Acetylcholinesterase) for neuromuscular junction function \parencite{rotundo_expression_2003}, \textbf{MEF2A} (Myocyte enhancer factor 2A) regulating differentiation \parencite{black_transcriptional_1998}, and \textbf{TRDN} (Triadin), a component of the calcium release complex \parencite{guo_association_1995}.

\subsubsection{Liver Developmental Markers}

Liver markers reflect metabolic maturation and circadian regulation.

Identified markers included \textbf{NPAS2} (Neuronal PAS domain protein 2) and \textbf{BHLHE41} (Basic helix-loop-helix family member e41), both circadian clock genes involved in the regulation of hepatic metabolism \parencite{zhang_night-restricted_2025}. \textbf{ALDH18A1} (Aldehyde dehydrogenase 18 family member A1), involved in proline biosynthesis \parencite{marco-marin_1_2020}, was also a top feature.

\subsubsection{Brain Developmental Markers}

Brain markers reflect neuronal maturation and epigenetic regulation.

Key markers were \textbf{DNMT3L} (DNA methyltransferase 3L), a regulator of de novo methylation \parencite{bourchis_dnmt3l_2001}, \textbf{HDAC6} (Histone deacetylase 6), involved in chromatin remodelling \parencite{hubbert_hdac6_2002}, and \textbf{KLK6} (Kallikrein-related peptidase 6), which is implicated in myelin biology \parencite{scarisbrick_kallikrein_2012}.

\subsubsection{Lung Developmental Markers}

Lung markers reflect immune development and RNA metabolism.

The top features included \textbf{IGF2BP2} (Insulin-like growth factor 2 mRNA-binding protein 2), a regulator of mRNA stability \parencite{li_hmga2-igf2bp2_2012}; \textbf{IGF2BP3}, a closely related oncofetal RNA-binding protein with a similar postnatal decline; \textbf{GFI1} (Growth factor independent 1), a transcription factor essential for haematopoiesis \parencite{van_der_meer_gfi1_2010}; and \textbf{NFIB} (Nuclear factor I B), known for its role in lung development \parencite{steele-perkins_transcription_2005}.

\subsubsection{Blood Developmental Markers}

Blood markers reflect haematopoietic development:

Markers included \textbf{GFI1}, essential for lymphocyte development \parencite{van_der_meer_gfi1_2010}, alongside other immune-related genes involved in lymphocyte differentiation and activation.

\subsection{Feature Selection Stability Analysis}

The stability of feature selection across random seeds was assessed to determine whether the same genes are consistently identified as important. Stability was quantified using the Jaccard similarity index between feature sets selected with different random seeds.

Mean Jaccard similarity ranged from 0.08 to 0.31 across tissues, evaluated across ten independent random seeds.

The moderate stability values reflect biological gene redundancy rather than model instability. Multiple gene sets can achieve similar classification performance because developmental changes involve coordinated regulation of many genes. Different random seeds may select different but functionally equivalent gene combinations.

Importantly, 6--42\% of top features appeared in $\geq$80\% of runs, indicating a core set of consistently selected markers surrounded by a variable periphery. This pattern is consistent with the presence of both essential developmental markers and redundant alternative markers.

\subsection{Expression-Stage Correlations}

The selected features were validated by assessing their correlation with developmental stage. For each tissue, the mean absolute Spearman correlation between top features and stage labels was computed.

Mean absolute Spearman correlations ranged from 0.2 to 0.6 across tissues, indicating that the selected features show substantial developmental regulation.

\subsection{Cross-Tissue Feature Overlap}

A striking finding was the minimal overlap of developmental markers across tissues. The pairwise Jaccard similarity between tissue-specific feature sets was only 0.2\%, indicating extremely low cross-tissue overlap.

This profound tissue-specificity indicates that each tissue uses distinct molecular programs for developmental progression. The genes that mark liver development are different from those that mark muscle or brain development, reflecting the tissue-specific nature of postnatal maturation.

The low cross-tissue overlap has important implications for staging applications. A single universal marker set cannot capture development across all tissues; instead, tissue-specific marker panels are required for accurate staging.

\subsection{Tissue-Specific Top Features}

To make the tissue-specific feature space directly inspectable, the genes in each tissue's ordinal model were ranked by aggregate LightGBM importance and the top~20 features per tissue were annotated via mygene.info and UniProt REST (\tabref{tab:S7}; row-$z$-scored expression trajectories in \figref{fig:S5}). For the four non-muscle tissues, biologically coherent stage-defining genes emerge that were not previously discussed in the cross-species muscle module. In brain, the ketogenic enzyme \textbf{HMGCS2} (infant peak) and the Hippo-pathway kinase \textbf{LATS1} (post-pubertal) appear in the top ranks, consistent with a metabolic-to-circuit-maturation transition during postnatal cortical development. In liver, \textbf{ALDH18A1} (proline/ornithine biosynthesis) and \textbf{COL1A1} peak in infancy alongside post-pubertal-peaking transcription factors \textbf{SATB2} and \textbf{CDH7}, mirroring the amino-acid-anabolism-to-mature-hepatocyte programme. In lung, the matricellular factor \textbf{CCN5} (post-pubertal; peak-vs-trough fold change 2.71) and the neural adhesion molecule \textbf{L1CAM} dominate alveolar maturation, whereas the oncofetal RNA-binder \textbf{IGF2BP2}---also a top cross-species muscle marker---signals neonatal mesenchymal growth. In blood, ribosomal \textbf{RPLP1} (infant peak; fold change 3.14) and the chemokine receptor \textbf{CCR6} (post-pubertal) track neonatal ribosomal/transcriptional surge and adaptive-immune maturation respectively. Placeholder rows are retained in \tabref{tab:S7} for top-ranked features that returned no curated description from either annotation source---predominantly novel pig loci and uncharacterised \texttt{LOC} entries---to flag these as opportunities for follow-up functional annotation.


\section{Functional Enrichment Analysis}\label{sec:results:enrichment}

\subsection{Tissue-Specific Functional Programs}

Functional enrichment analysis of stage-informative genes revealed distinct developmental programs in each tissue. The results are summarised in \figref{fig:enrichment} and described below.

\begin{figure}[H]
\centering
\includegraphics[width=\textwidth]{fig3_functional_enrichment.pdf}
\caption{\textbf{Functional enrichment analysis of stage-informative genes.}
Functional network showing interconnected enrichment modules, with nodes representing enriched terms coloured by tissue and edges connecting terms sharing >25\% of genes (Jaccard similarity). Hubs and cluster representatives are labelled with biological process or pathway names.}
\label{fig:enrichment}
\end{figure}

\subsubsection{Brain: Synaptic Refinement and Metabolic Support}

Brain developmental markers were enriched for processes related to synaptic function and neuronal metabolism. Enrichment for sarcosine-related processes reflects the role of this glycine derivative as a co-agonist at NMDA receptors, which is critical for synaptic plasticity \parencite{zhang_glycine_2009}. Furthermore, the enrichment of beta-tubulin binding reflects the importance of microtubule dynamics for axonal transport and neuronal morphology.

These findings align with the known prolonged development of neural circuits during postnatal life, including ongoing synaptogenesis, synaptic pruning, and myelination \parencite{jiang_cellular_2016, kang_spatio-temporal_2011, semple_brain_2013}.

\subsubsection{Liver: Metabolic Autonomy and Homeostasis}

Liver markers were enriched for metabolic processes. Enrichment for autophagosome-membrane adaptors reflects the role of autophagy in hepatic homeostasis and the response to nutritional changes associated with weaning \parencite{gordillo_orchestrating_2015, pluske_gastrointestinal_2018}. There was also a general enrichment for metabolic enzyme expression, reflecting the liver's central role in nutrient processing.

The functional profile reflects the liver's adaptation to postnatal metabolic demands, including the switch from maternal nutrition to independent feeding.

\subsubsection{Muscle: Contractile Machinery and Neuromuscular Signalling}

Muscle markers were enriched for contractile machinery and neuromuscular signalling processes. Enrichment for acetylcholinesterase activity \parencite{rotundo_expression_2003} reflects the maturation of neuromuscular junctions and the establishment of efficient synaptic transmission. Additionally, the enrichment for sarcomere organisation captures the development of the mature contractile apparatus.

These findings capture the establishment of adult muscle function, including the transition from neonatal to adult fibre types \parencite{lefaucheur_pattern_1995, lefaucheur_second_2010} and the refinement of neuromuscular connectivity.

\subsubsection{Lung: Immune Establishment and RNA Metabolism}

Lung markers were enriched for immune establishment and RNA metabolism processes. Enrichment for IL-6 signalling regulation reflects the establishment of pulmonary immune function and inflammatory responses. Concurrently, enrichment for N6-methyladenosine (m6A) modification \parencite{wang_n6-methyladenosine_2014} suggests a role for post-transcriptional regulation in lung development.

The lung's exposure to environmental antigens drives rapid immune maturation during postnatal life \parencite{georgountzou_postnatal_2017}, as captured by these functional enrichments.

\subsection{Integration Across Tissues}

The functional network in \figref{fig:enrichment} reveals both tissue-specific and shared developmental processes. While the majority of enriched terms are specific to individual tissues, some processes are shared across tissues, such as cell proliferation and growth regulation (muscle, liver, lung), immune function (blood, lung), and metabolic adaptation (liver, muscle). These shared processes reflect the common challenges of postnatal development, including the need for growth, immune protection, and metabolic adaptation across all tissues.


\section{Cross-Species Validation}\label{sec:results:crossspecies}

\subsection{Conservation of Developmental Trajectories}

Cross-species analysis with human skeletal muscle transcriptomes \parencite{schaiter_molecular_2024} ($n = 11$) provided evidence of conservation of developmental gene expression programs. The results are summarised in \figref{fig:crossspecies}.

\begin{figure}[H]
\centering
\includegraphics[width=\textwidth]{fig4_cross_species.pdf}
\caption{\textbf{Cross-species validation of developmental gene expression.}
Human skeletal muscle transcriptomic data comparing infant (7--28 months) versus adult (30--56 years) samples were used to validate conservation of developmental signatures.
\textbf{(a)} Scatter plot of log$_2$ fold changes between pig and human, showing significant positive correlation.
\textbf{(b)} Feature importance of top conserved markers in the developmental stage classifier.
\textbf{(c)} Paired fold change comparison showing pig and human values for each conserved gene.}
\label{fig:crossspecies}
\end{figure}

\subsubsection{Correlation Between Species}

Using pre-specified significance and effect-size criteria (FDR $< 0.10$, Benjamini--Hochberg corrected, and $|\log_2\mathrm{FC}| > 0.5$ in both species), 66 genes with significant and biologically meaningful developmental fold-changes were identified. The Pearson correlation between pig and human fold changes was $r = 0.693$ ($p = 1.09 \times 10^{-10}$, 95\% CI: 0.542--0.801), indicating a strong positive relationship between developmental trajectories in the two species.

Bootstrap analysis confirmed the robustness of this correlation. The 95\% bootstrap confidence interval was 0.50--0.82 (based on 1,000 iterations), excluding zero and confirming the statistical significance of the conservation signal.

\subsubsection{Directional Concordance}

Of the 66 genes meeting significance criteria, 64 (97.0\%) showed the same direction of change in both species. This high directional concordance indicates that genes upregulated during pig muscle development tend also to be upregulated during human muscle development, and vice versa.

The 2 genes showing discordant directions (3.0\%) may reflect species-specific differences in intramuscular lipid metabolism or sampling variability.

\subsection{Sensitivity Analysis}

The robustness of the cross-species correlation was assessed across a range of FDR thresholds. Sensitivity analysis across thresholds from 0.05 to 0.20 showed correlations ranging from $r = 0.61$ to $r = 0.69$, with the correlation magnitude decreasing slightly at more permissive thresholds as noisier genes are included. This consistency demonstrates that the conservation signal is not an artefact of the specific parameter choices.

The sensitivity analysis results are presented in \figref{fig:S1}, showing the correlation coefficient as a function of FDR threshold.

\subsection{Conserved Functional Modules}

Three functional modules of conserved developmental genes were identified:

\subsubsection{Module 1: Neuromuscular Maturation and Growth Signalling}

This module comprises genes with high infant expression that decrease during postnatal development, reflecting the transition from active myogenesis to mature muscle homeostasis:

\textbf{ACHE} (pig FC = $-$3.34, human FC = $-$0.53), essential for neuromuscular junction function \parencite{rotundo_expression_2003}, shows elevated expression in the infant stage (0--20 days; mean TPM = 17.5) that progressively declines through early childhood (TPM = 7.6), pre-pubertal (TPM = 3.4), to baseline in post-pubertal stages (TPM = 1.3). \textbf{NRK} (pig FC = $-$8.56, human FC = $-$2.45), an embryonic Ste20-type kinase expressed predominantly in developing skeletal musculature \parencite{kanai-azuma_nrk_1999}, has the second-largest fold change, with infant TPM = 128.2 declining to 0.3 in adults. \textbf{KREMEN1} (pig FC = $-$1.69, human FC = $-$1.43), a negative regulator of Wnt/$\beta$-catenin signalling \parencite{osada_wnt_2006, qin_canonical_2024}, shows a similar gradual postnatal decline. \textbf{IGF2BP2} (pig FC = $-$2.19, human FC = $-$2.03) and \textbf{IGF2BP3} (pig FC = $-$4.85, human FC = $-$1.23), oncofetal RNA-binding proteins that regulate myoblast proliferation via translational control of growth-related transcripts \parencite{li_hmga2-igf2bp2_2012}, show pronounced postnatal decline. \textbf{SORCS2} (pig FC = $-$4.67, human FC = $-$0.95), a regulator of motor neuron innervation \parencite{thomasen_sorcs2_2023}, is highly expressed during early postnatal life (infant TPM = 14.4) and becomes progressively downregulated through early childhood (TPM = 3.0), pre-pubertal (TPM = 2.2), reaching its lowest in post-pubertal (TPM = 0.6). \textbf{ETV5} (pig FC = $-$2.80, human FC = $-$2.24), a member of the PEA3 subfamily of ETS transcription factors, which has been implicated in myogenic differentiation \parencite{taylor_role_1997}, and \textbf{CCDC8} (pig FC = $-$4.44, human FC = $-$3.14), a core component of the 3M growth-regulatory complex \parencite{hanson_exome_2011}, complete this module.

\subsubsection{Module 2: Structural Refinement and ECM Remodelling}

This module reflects the transition from rapid tissue growth to structural consolidation, encompassing both cytoskeletal dynamics and extracellular matrix maturation:

\textbf{TCAF1} (pig FC = $-$4.72, human FC = $-$1.02), which regulates TRPM8 ion channel trafficking and calcium signalling \parencite{gkika_trp_2015}, is the highest-importance gene among the cross-species markers, with infant TPM = 78.9 declining to 3.0 in adults. \textbf{FGFRL1} (pig FC = $-$3.03, human FC = $-$2.10) is required for slow muscle fibre development \parencite{amann_fgfrl1_2014}; its expression peaks during the infant stage (TPM = 77.4), drops sharply in early childhood (TPM = 15.6), and decreases to its lowest in post-pubertal (TPM = 4.3). \textbf{DLK1} (pig FC = $-$9.15, human FC = $-$2.88), a crucial regulator of the myogenic program \parencite{waddell_dlk1_2010,andersen_dual_2013}, shows high infant expression (TPM = 338.8) declining sharply as the muscle fibre population matures. \textbf{POSTN} (pig FC = $-$7.97, human FC = $-$2.93), a secreted matricellular protein associated with ECM remodelling during myofibre maturation \parencite{ozdemir_periostin_2014}, is the most highly expressed ECM gene in this set (infant TPM = 602.4), declining dramatically as tissue architecture stabilises. \textbf{WDR1} (pig FC = $-$1.69, human FC = $-$0.52), a cofactor of cofilin-mediated actin filament disassembly \parencite{ono_functions_2018}, \textbf{FSCN1} (pig FC = $-$4.28, human FC = $-$1.95), an actin-bundling protein driving cell migration during morphogenesis \parencite{adams_roles_2004}, \textbf{FREM2} (pig FC = $-$4.58, human FC = $-$1.05), a basement membrane protein \parencite{kiyozumi_breakdown_2006}, \textbf{ADAMTS2} (pig FC = $-$4.25, human FC = $-$1.83), the principal procollagen N-proteinase \parencite{colige_cdna_1997}, and \textbf{COL11A2} (pig FC = $-$8.35, human FC = $-$0.70), a collagen component regulating fibril organisation \parencite{sun_collagen_2020}, reflect intensive ECM biosynthesis during rapid postnatal growth. \textbf{TRIM54}, a muscle-specific E3 ubiquitin ligase involved in sarcomere assembly \parencite{gregorio_functional_2005}, shows the opposite pattern---lower infant expression (TPM = 78.3) increasing sharply in early childhood (TPM = 351.1) as the mature contractile apparatus is established.

\subsubsection{Module 3: Metabolic Transition}

This module captures genes upregulated during metabolic maturation of skeletal muscle:

\textbf{ME1} (pig FC = 2.59, human FC = 0.64) encodes malic enzyme 1 \parencite{simmen_malic_2020}, a key source of NADPH for lipid biosynthesis; its expression increases progressively from infant (TPM = 2.8) to adult stages (TPM = 16.9), reflecting the transition from carbohydrate to lipid metabolism. \textbf{MAOB} (pig FC = 2.55, human FC = 0.80), a mitochondrial outer membrane enzyme catalysing oxidative deamination of monoamines, shows concordant postnatal upregulation (infant TPM = 5.7 to adult TPM = 33.4), consistent with the general increase in mitochondrial content during postnatal muscle maturation. Notably, \textbf{PLA2G4E} (pig FC = 2.32, human FC = $-$1.25), a group~IVE cytosolic phospholipase, is the only discordant gene among the top 20 markers, showing opposite developmental trajectories in pig and human---potentially reflecting species-specific differences in intramuscular lipid metabolism.

\subsection{Biological Significance of Conservation}

The conservation of developmental trajectories between pig and human muscle has important implications for the use of pigs as translational models. The shared expression changes indicate that:

First, the porcine developmental staging framework captures evolutionarily conserved processes rather than species-specific artefacts. Second, insights from pig developmental biology are likely to translate to human biology, enhancing the pig's value as a biomedical model. Finally, the identified developmental markers likely represent fundamental regulators maintained across 80 million years of mammalian evolution \parencite{kumar_timetree_2022}.

\subsection{Functional Annotations of Key Markers}

The biological functions of key conserved markers provide insight into the molecular basis of postnatal muscle development. Selected annotations are provided in \tabref{tab:annotations}.

\begin{table}[H]
\centering
\caption{\textbf{Functional annotations of key conserved developmental markers.} FC = $\log_2$ fold change (adult/infant).}
\label{tab:annotations}
\footnotesize
\begin{tabular}{lp{5.5cm}cc}
\toprule
\textbf{Gene} & \textbf{Developmental Function} & \textbf{Pig FC} & \textbf{Human FC} \\
\midrule
\textit{SORCS2} & VPS10 domain-containing receptor that binds progranulin to regulate motor neuron development. & $-$4.67 & $-$0.95 \\
\addlinespace
\textit{FGFRL1} & Required for slow muscle fibre development; loss of function leads to diaphragm defects. & $-$3.03 & $-$2.10 \\
\addlinespace
\textit{KREMEN1} & Negative regulator of Wnt/$\beta$-catenin signalling, controlling cell survival and tissue patterning. & $-$1.69 & $-$1.43 \\
\addlinespace
\textit{ACHE} & Key enzyme at the neuromuscular junction and structural adhesion molecule guiding neuronal growth. & $-$3.34 & $-$0.53 \\
\addlinespace
\textit{DLK1} & Regulator of the myogenic program; inhibits myoblast proliferation while promoting differentiation. & $-$9.15 & $-$2.88 \\
\addlinespace
\textit{TCAF1} & Regulates TRPM8 ion channel trafficking and Ca$^{2+}$ signalling during muscle growth. & $-$4.72 & $-$1.02 \\
\addlinespace
\textit{NRK} & Embryonic Ste20-type kinase promoting actin polymerisation during skeletal muscle formation. & $-$8.56 & $-$2.45 \\
\bottomrule
\end{tabular}
\end{table}

\subsection{Limitations of Cross-Species Analysis}

Several limitations of the cross-species analysis must be acknowledged:

First, the human dataset was small ($n = 11$), limiting statistical power. Second, the substantial age gap between human infant and adult groups prevented characterisation of intermediate stages. Third, validation was limited to muscle tissue due to data availability. Fourth, although the gene selection criteria (FDR $< 0.10$, $|\log_2\mathrm{FC}| > 0.5$) were pre-specified, sensitivity analysis confirms robustness across a range of thresholds.

Despite these limitations, the conservation analysis provides important validation of the biological accuracy of the porcine developmental staging framework.


\section{Within-Project Batch Validation}\label{sec:results:batch}

To assess whether cross-project technical variation confounds the developmental signals, we performed a within-project fold-change validation. We identified the two largest single-project muscle cohorts that each span multiple developmental stages under uniform protocols: PRJNA488311 ($n_{\text{total}} = 81$) and PRJNA486202 ($n_{\text{total}} = 75$), each containing samples from infant through post-pubertal stages.

Within-project fold-changes correlated strongly with full-dataset estimates across three developmental contrasts (infant vs.\ early childhood, infant vs.\ pre-pubertal, and infant vs.\ post-pubertal), with pooled Pearson $r = 0.67$--$0.81$ for all expressed genes (\tabref{tab:S5}). Critically, the 66 cross-species validated genes showed even higher agreement ($r = 0.76$--$0.98$), confirming that the developmental signatures reflect biology rather than batch artefacts even for the earliest and subtlest developmental transitions.

These results provide important evidence that the multi-project aggregation strategy does not introduce spurious developmental signals, and that the key findings---particularly the cross-species conserved genes---are robust to technical variation across sequencing platforms and library protocols.
